\documentclass[conference]{IEEEtran}
% \IEEEoverridecommandlockouts
% The preceding line is only needed to identify funding in the first footnote. If that is unneeded, please comment it out.
\usepackage{cite}
\usepackage{amsmath,amssymb,amsfonts}
\usepackage{algorithmic}
\usepackage{graphicx}
\usepackage{textcomp}
\usepackage{xcolor}
\def\BibTeX{{\rm B\kern-.05em{\sc i\kern-.025em b}\kern-.08em
    T\kern-.1667em\lower.7ex\hbox{E}\kern-.125emX}}
\begin{document}

\title{Support Edge Based Classifiers Review}

\author{\IEEEauthorblockN{Eduardo Henrique Basilio de Carvalho}
\IEEEauthorblockA{\textit{Departamento de Engenharia Eletrônica} \\
\textit{Universidade Federal de Minas Gerais}\\
Belo Horizonte, Brasil \\
eduardohbc@ufmg.br}
}

\maketitle

\begin{abstract}
%
% In this paper, we:
%   - review the theoretical background of Support Edge based classifiers;
%   - evaluate its implementation in a modern compiled programming language;
%   - evaluate possible improvements in the implementation;
%   - compare a hardware implementation achieved by high-level synthesis with state-of-the-art hardware description language implementations;
\end{abstract}

\begin{IEEEkeywords}
Classification, Machine Learning, Nearest Neighbor, Gabriel Graph, Support Edge
\end{IEEEkeywords}

\section{Introduction}
%
%
% Talk about each of the references and the room for improvement
%
%

\section{Theoretical Background}
%
%
% Provide theoretical background on the following topics:
%   - Gabriel Graph
%   - Support Edge
%   - CHIP-CLAS
%   - RCHIP-CLAS
%   - NN-CLAS
%
%

\section{Possible Improvements}
%
% Possible improvements in the implementation:
%   - use of adjacency list instead of adjacency matrix
%   - after filtering, only compute the edge of points of opposite classes
%   

\section{Methodology}
%
%
% Describe the methodology used to evaluate the implementation:
%   - implementation can be accessed in https://github.com/eduardo-ufmg/CLAS
%   - the classifiers are writen in C++20
%   - currently, synthetic data is used to evaluate the classifiers
%   - the data is generated with python3 sklearn
%   - evaluation scripts are also written in python3 and available in the same repository
%   - HLS implementation is done using Vitis HLS for two different platforms:
%     - U55C
%     - Zedboard
%

\section{Results}
%
% figure: blob_good
% figure: blob_bad
% figure: spiral
% figure: moons
% figure: circles
% figure: xor
%

\section*{References}

\begin{thebibliography}{00}
\bibitem{torres2016} L. C. B. Torres, "Classificador por arestas de suporte (CLAS): métodos de aprendizado baseados em Grafos de Gabriel," Manuscript, 2016.
\bibitem{souza2019} A. C. Souza, C. Leite Castro, J. A. Garcia, L. C. B. Torres, L. J. Acevedo Jaimes and B. R. A. Jaimes, "Improving the Efficiency of Gabriel Graph-based Classifiers for Hardware-optimized Implementations," 2019 XXII Symposium on Image, Signal Processing and Artificial Vision (STSIVA), Bucaramanga, Colombia, 2019.
\bibitem{arias2021} J. Arias-Garcia et al., "Enhancing Performance of Gabriel Graph-Based Classifiers by a Hardware Co-Processor for Embedded System Applications," in IEEE Transactions on Industrial Informatics, vol. 17, no. 2, Feb. 2021.
\bibitem{arias2024} J. Arias-Garcia et al., "Improved Design for Hardware Implementation of Graph-Based Large Margin Classifiers for Embedded Edge Computing," in IEEE Transactions on Neural Networks and Learning Systems, vol. 35, no. 1, Jan. 2024.
\bibitem{torres2015} L. C. B. Torres, C. L. Castro and A. P. Braga, "A parameterless mixture model for large margin classification," 2015 International Joint Conference on Neural Networks (IJCNN), Killarney, Ireland, 2015.
\bibitem{torres2021} L. C. B. Torres, C. L. Castro, F. Coelho and A. P. Braga, "Large Margin Gaussian Mixture Classifier With a Gabriel Graph Geometric Representation of Data Set Structure," in IEEE Transactions on Neural Networks and Learning Systems, vol. 32, no. 3, March 2021.
\bibitem{torres2015b} L. C. B. Torres, C. L. Castro, F. Coelho, F. Sill Torres and A. P. Braga, "Distance-based large margin classifier suitable for integrated circuit implementation," Manuscript, 2015.
\end{thebibliography}

\end{document}
